% mazzi: 10
\chapter{L'elettromiografia di superficie}

\textbf{Elettromiografia}: disciplina sperimentale legata alla generazione, alla misura, all'analisi
e all'utilizzo del segnale elettrico emanato dai muscoli; l'obiettivo è valutare il funzionamento
dei muscoli attraverso i potenziali elettrici che essi generano.

La misura \textbf{EMGs} sincronizzata con la dinamometria isoinerziale è usata nella valutazione
funzionale sportiva per studiare l'\textbf{attivazione neuromuscolare} durante l'esecuzione di
movimenti funzionali e specifici: la dinamometria scandisce le fasi del movimento e permette così di
collocare ogni segnale nella fase in cui è stato prodotto.

% ---------------------------------------------------------------------------
\section{Il substrato del segnale}

\rubrica{L'unità motoria}
\textbf{Unità motoria} (MU): la più piccola unità funzionale che descrive il controllo neurale della
contrazione muscolare, composta dal \textbf{motoneurone $\alpha$}, dalle diramazioni del suo assone
e dalle fibre muscolari innervate.
\begin{itemize}
  \item il termine \textit{unità} indica che tutte le fibre della singola MU formano un unico blocco
    funzionale con l'assone che le innerva e si contraggono tutte insieme in seguito al segnale
    proveniente dall'assone;
  \item un singolo assone fa contrarre in contemporanea un numero di fibre pari a quelle della sua
    MU;
  \item il numero di MU varia molto da muscolo a muscolo: circa 100 nei piccoli muscoli della mano,
    oltre 1000 nei muscoli più grandi degli arti superiori o inferiori.
\end{itemize}

\rubrica{Reclutamento e frequenza di attivazione}
Durante le contrazioni volontarie la forza esercitata è modulata da due parametri indipendenti tra
loro:
\begin{enumerate}
  \item \textbf{reclutamento delle MU}: all'aumentare delle MU attivate aumenta la forza;
  \item \textbf{frequenza di attivazione} delle MU (\textit{firing}): all'aumentare della frequenza
    aumenta la forza --- è il meccanismo con cui la tensione cresce dalla forza dinamica massima
    alla forza isometrica massima, senza reclutare nuove fibre.
\end{enumerate}
La combinazione dei due parametri varia da muscolo a muscolo, varia con la velocità del movimento e
dipende dalla \textbf{fatica muscolare}: aspetto molto complesso e ancora dibattuto
scientificamente.

\rubrica{Eccitabilità della membrana e propagazione}
\begin{enumerate}
  \item dalla \textbf{giunzione neuromuscolare} (placca motrice) il potenziale d'azione si sposta
    sulla fibra in entrambe le direzioni, dalla placca verso i tendini;
  \item l'eccitazione causa il rilascio di ioni \textbf{calcio} ($Ca^{2+}$) nello spazio
    intracellulare e la conseguente contrazione della fibra, per una serie di processi chimici
    concatenati: è l'\textbf{accoppiamento elettromeccanico};
  \item la propagazione si basa sulla generazione di nuovi potenziali d'azione nei punti successivi
    della fibra: l'insorgenza in un punto crea una differenza di potenziale con le zone vicine,
    ancora a riposo;
  \item tra zona attiva e zona inattiva nasce una \textbf{corrente locale}, che depolarizza la zona
    inattiva fino alla soglia $\rightarrow$ nuovo potenziale d'azione, e così via.
\end{enumerate}

% ---------------------------------------------------------------------------
\section{La registrazione del segnale}

Il segnale si capta con elettrodi applicati sulla pelle: elettrodi $\rightarrow$ amplificatore
$\rightarrow$ computer, che lo elabora.

\begin{itemize}
  \item tutto ciò che si interpone tra muscolo ed elettrodi altera il segnale, in particolare il
    \textbf{grasso} del tessuto sottocutaneo;
  \item per ottenere un segnale netto e privo di rumore di fondo la superficie va pulita
    accuratamente.
\end{itemize}

\rubrica{Applicazione degli elettrodi}
\begin{enumerate}
  \item \textbf{depilazione} della zona (nell'esempio il vasto laterale);
  \item \textbf{scrub} con una carta vetrata finissima, per eliminare cellule morte e impurità della
    superficie;
  \item due \textbf{elettrodi affiancati}, a distanza ravvicinata di massimo 1,5\,cm, con i cavi
    applicati direttamente;
  \item un \textbf{terzo elettrodo} in zona neutra, che serve a calcolare il differenziale rispetto
    a quelli posizionati sul muscolo;
  \item l'applicazione segue l'\textbf{angolo di pennazione} del muscolo, cioè la direzione delle
    fibre.
\end{enumerate}
Esistono \textbf{tavole anatomiche} che indicano i siti per gli elettrodi a filo (\textit{fine wire
sites}) e per quelli di superficie (\textit{surface sites}), utili a individuare le zone di
innervazione dei muscoli superficiali.

\rubrica{Posizionamento rispetto alla zona di innervazione}
L'elettrodo va posto lontano dalla zona di innervazione, perché da lì la depolarizzazione si propaga
in entrambe le direzioni, a monte e a valle:
\begin{itemize}
  \item \textbf{attorno alla zona di innervazione} il segnale non è ben definito;
  \item esiste una zona in cui il segnale è molto più forte;
  \item il segnale degrada man mano che il muscolo si trasforma in tendine $\rightarrow$ nullo sul
    tendine.
\end{itemize}

\rubrica{Che cosa si misura realmente}
La \textbf{somma algebrica} di tutti i potenziali d'azione trasmessi sulla fibra muscolare in quel
preciso istante. Il segnale viene poi filtrato, amplificato e riprodotto in forma grafica dal
computer per l'analisi; quello registrato senza elaborazione è il \textbf{segnale grezzo}.

% ---------------------------------------------------------------------------
\section{Le analisi possibili}

\rubrica{Analisi qualitativa}
\begin{itemize}
  \item \textbf{attivazione o deattivazione} muscolare dei muscoli su cui sono applicati gli
    elettrodi;
  \item \textbf{fasi dell'attivazione} in un movimento specifico, cioè come avviene la
    depolarizzazione nel tempo;
  \item \textbf{verifica dell'apprendimento} di particolari movimenti: in registrazioni successive
    si osserva se l'attivazione resta la stessa o si abbassa --- in un confronto pre/post
    l'elettrodo deve stare esattamente nella stessa posizione;
  \item \textbf{ampiezza del segnale};
  \item \textbf{numero di \textit{zero crossing}}: i passaggi per lo zero del segnale grezzo;
  \item \textbf{numero e/o ampiezza degli \textit{spike}}, cioè dei picchi.
\end{itemize}

\rubrica{Analisi quantitativa}
\begin{itemize}
  \item \textbf{voltaggio};
  \item \textbf{integrazione}: calcolo dell'area sottesa, come sommatoria di tutte le piccole aree
    sottese dai singoli \textit{spike};
  \item \textbf{RMS} (\textit{root mean square}, scarto quadratico medio).
\end{itemize}

\subsection{Elaborazione del segnale: ampiezza e \textit{smoothing}}

Il \textbf{pattern di interferenza} del segnale EMG è di natura casuale e non riproducibile: lo
stesso task motorio genera scariche di impulsi simili ma mai uguali.
\begin{itemize}
  \item si cerca perciò di eliminare la parte non riproducibile con algoritmi di
    \textbf{\textit{smoothing}}, che evidenziano l'andamento del trend medio;
  \item l'RMS, elevando al quadrato, rende positivo l'intero segnale $\rightarrow$ il segno si
    elimina;
  \item l'\textbf{iEMG} (EMG integrato) è il risultato della sommatoria delle aree sottese dai
    singoli \textit{spike}: una curva crescente che si azzera a ogni \textit{reset}.
\end{itemize}

% ---------------------------------------------------------------------------
\section{Valutazione neuromuscolare applicata}

\rubrica{Allenamento pliometrico e con sovraccarichi (Hakkinen, 1985)}
Con l'iEMG si confrontano gli adattamenti di due tipi di allenamento, misurati su forza isometrica e
attività elettromiografica:
\begin{itemize}
  \item \textbf{pliometrico}: PF $+11\,\%$, Max.\ RFD $+24\,\%$, iEMG $+38\,\%$ nella fase iniziale
    e $+8\,\%$ a regime;
  \item \textbf{sovraccarichi}: PF $+27\,\%$, Max.\ RFD $+0,4\,\%$, iEMG $+3\,\%$.
\end{itemize}
L'aumento di forza è dunque maggiore con i sovraccarichi, mentre il guadagno di attività
elettromiografica è nettamente maggiore con il pliometrico $\rightarrow$ l'adattamento del
pliometrico è prevalentemente \textbf{neurale}, legato alla \textbf{sincronizzazione} tra muscoli
agonisti e antagonisti.

\subsection{Il comportamento neuromuscolare sulle superfici instabili}

L'ipotesi di partenza era che sollevare carichi su una superficie instabile, richiedendo uno sforzo
maggiore, permettesse di ridurre il carico sulle spalle a parità di stimolo: prima di darlo per
acquisito, il comportamento neuromuscolare è stato analizzato.

\rubrica{Protocollo}
\begin{itemize}
  \item riferimento: \textbf{squat parallelo} su superficie stabile;
  \item stessa esecuzione su instabilità crescente: \textit{medusa} con appoggio sulla superficie
    piana, poi \textit{medusa} con appoggio convesso (rovesciata);
  \item i tre esercizi ripetuti anche su \textbf{superficie vibrante} (NEMES, BoscoSystem) a 30\,Hz,
    e nelle combinazioni vibrazione + medusa e vibrazione + medusa rovesciata.
\end{itemize}

\rubrica{Risultati}
\begin{itemize}
  \item \textbf{EMGs \% rispetto alla superficie stabile}: al crescere dell'instabilità l'attività
    del \textbf{quadricipite} --- muscolo \textbf{dinamico}, il motore che permette di sviluppare
    alti gradienti di forza --- scende sotto il valore di riferimento, mentre quella del
    \textbf{tibiale} --- muscolo \textbf{posturale} --- aumenta di parecchie volte, con il massimo
    nella condizione vibrazione + medusa rovesciata;
  \item \textbf{EMG totale e tempo di esecuzione}: crescono entrambi con l'instabilità.
\end{itemize}

\rubrica{Interpretazione}
\begin{itemize}
  \item il sistema neuromuscolare adotta una strategia per rimanere in equilibrio, a scapito dello
    sviluppo di forza;
  \item tempi di esposizione eccessivamente lunghi a questo tipo di lavoro condizionano la
    muscolatura posturale più di quella dinamica, con decremento dei gradienti di forza;
  \item serve uno sforzo enorme del sistema neuromuscolare per produrre un movimento lento:
    l'opposto dell'obiettivo prestativo, che è ridurre i tempi e aumentare la velocità di
    contrazione.
\end{itemize}

% ---------------------------------------------------------------------------
\section{Lettura dei tracciati sincronizzati}

Negli esempi che seguono sono sincronizzati sullo stesso tempo l'\textbf{EMGrms} dei muscoli
indagati --- vasto laterale e vasto mediale destri come agonisti, \textit{hamstring} destro come
antagonista --- e velocità e posizione del carico.

\rubrica{Ritardo elettromeccanico}
\textbf{Ritardo elettromeccanico}: tempo che intercorre tra l'arrivo del segnale al muscolo --- cioè
l'inizio della depolarizzazione della membrana, il primo istante che l'EMG può registrare --- e
l'inizio effettivo del movimento; dura normalmente circa 100\,ms.
\begin{itemize}
  \item in quell'intervallo avvengono la generazione della tensione muscolare e la sua trasmissione
    agli elementi elastici in serie, al tendine e all'osso;
  \item il movimento inizia solo quando la tensione supera la resistenza esterna;
  \item ridurre questo tempo è importante per la prestazione, soprattutto nel reagire velocemente.
\end{itemize}

\subsection{Squat con 30\,kg, modalità concentrica}
\begin{itemize}
  \item l'\textbf{attivazione} precede l'inizio del movimento, riconoscibile dal punto in cui
    cambiano posizione e velocità;
  \item il \textbf{picco} di attivazione degli agonisti si colloca all'inizio del movimento;
  \item poi l'attivazione degli agonisti diminuisce e, per \textbf{innervazione reciproca}, aumenta
    quella dell'antagonista $\rightarrow$ frena il movimento e riporta a zero la velocità.
\end{itemize}

\subsection{Squat con 30\,kg, eccentrico-concentrico}
\begin{enumerate}
  \item nella \textbf{discesa} l'attività elettrica è nulla: il sistema asseconda l'accelerazione di
    gravità;
  \item da un certo punto il muscolo genera una tensione che lo fa contrarre mentre si allunga,
    accumulando \textbf{energia elastica} negli elementi elastici in serie e in parallelo;
  \item il picco si osserva nel \textbf{punto di inversione}, quando la velocità torna a zero e poi
    diventa positiva: è l'inizio della fase propulsiva, ed è lì che conviene lavorare per migliorare
    la spinta;
  \item il \textbf{tempo di accoppiamento} tra le due fasi deve restare breve: se si allunga
    $\rightarrow$ l'energia elastica e il riflesso da stiramento accumulati si perdono.
\end{enumerate}

\subsection{Squat jump}
\begin{itemize}
  \item essendo un movimento concentrico, l'attivazione massima si ha subito dopo l'inizio,
    leggermente spostata in avanti rispetto agli esempi precedenti; segue la \textbf{fase di volo};
  \item alle velocità superiori l'attivazione dell'\textbf{\textit{hamstring}} diventa più
    importante di quanto osservato con il carico di 30\,kg: comincia nella fase concentrica, per
    innervazione reciproca dopo che gli agonisti hanno esaurito la loro azione;
  \item quell'attivazione serve anche a estendere il tronco sugli arti inferiori.
\end{itemize}
Va ricordato che movimenti puramente concentrici non esistono in natura: tutti i movimenti naturali
contro gravità usano la strategia del ciclo stiramento-accorciamento.

\subsection{Counter movement jump}
\begin{itemize}
  \item nella \textbf{fase eccentrica} l'attivazione dell'\textit{hamstring} è bassa;
  \item nella fase successiva il comportamento è quello del movimento concentrico, con il picco
    nella fase di volo;
  \item al punto di impatto della ricaduta si osserva una \textbf{preattivazione}: il sistema
    neuromuscolare attiva la muscolatura prima del contatto con il terreno, preparandosi
    all'impatto.
\end{itemize}

\subsection{Movimento esplosivo a catena cinetica aperta}
Flesso-estensione della gamba sulla coscia, senza fase di volo:
\begin{enumerate}
  \item nella prima fase il soggetto flette il ginocchio verso l'alto: si attiva il
    \textbf{flessore}, per mantenere la posizione di circa $90^\circ$ della gamba sulla coscia;
  \item nell'\textbf{estensione} l'attivazione dei due vasti è sincrona;
  \item essendo il movimento velocissimo, l'\textit{hamstring} si attiva verso la fine;
  \item \textbf{picco di attivazione degli agonisti} $\rightarrow$ inibizione degli antagonisti
    (innervazione reciproca); passato quel periodo l'azione degli antagonisti riprende.
\end{enumerate}

\rubrica{Il timing agonisti/antagonisti}
Sia nella \textbf{ricaduta} sia nella \textbf{fase propulsiva} l'attività di agonisti e antagonisti
deve avere un timing perfetto: una variazione nel timing di attivazione --- anche di un millesimo di
secondo di anticipo --- può creare problemi al muscolo più debole, cioè l'antagonista, fino
all'\textbf{infortunio dei flessori}.

\rubrica{Perché serve l'EMGs}
L'elettromiografia di superficie fornisce informazioni più dettagliate sul comportamento del sistema
neuromuscolare durante la prestazione, che il solo parametro biomeccanico non può dare: permette di
suddividere le fasi del movimento e di capire attivazione, disattivazione e \textbf{rapporto
agonisti/antagonisti} in ciascuna fase.
